\documentclass[../main.tex]{subfiles}

\begin{document}

\section{Conclusion}
\label{sec:conclusion}

\subsection{Overall Assessment Conclusion}

Within the engagement scope and the August 27, 2026 evidence cutoff, this
assessment supports a bounded conclusion: HDF5 input is not necessarily inert
data. File-controlled metadata can reach memory, allocation, traversal,
external-resource, plugin, deserialization, storage, and disclosure boundaries.
Risk arises at the component that supplies and controls the relevant authority,
and not every appearance of HDF5 in an adverse event is therefore a defect in
the core library or file format.

The evidence supports four ecosystem-level judgments. First, historically
affected core and official-tool paths exhibit recurring failures to validate
file-derived values before trusted C operations; local fixes can treat a known
variant, but class elimination requires invariant-driven prevention and
continuing regression evidence. Second, applications can turn documented HDF5
capabilities into security mechanisms when they process untrusted files with
ambient filesystem, plugin-loading, deserialization, or resource authority.
Third, reader acceptance, format conformance, and safe deployment are separate
decisions: permissive historical acceptance is not proof that a file is
consistent, and stricter rejection is not automatically a regression --- and in
at least one documented case the reference library, not the independent writer,
was the non-conforming party, while in another it is still undetermined which
artifact is wrong. Fourth,
privacy harm can arise during correct and authorized operation because payload
protection does not necessarily protect metadata, structure, derived
information, caches, logs, or other workflow artifacts. These four judgments do
not exhaust the safety surface: storage exhaustion, abrupt termination, and
filter-dependent loss of interpretability also threaten data that no adversary
ever touched. This revision registers those bounded scenarios while leaving
unmeasured VFD, VOL, and specification-governance questions in the evidence
backlog.

\begin{notebox}
The registered scenarios have a determinate, conditional mitigation-priority
order, and confirmed applicable scope follows it. No exception permits a known,
uncontained Critical-equivalent external-resource path to process untrusted
input; any other variance requires a scoped, accountable, expiring record and
tested interim control. The rubric and response routing control---not chapter
order, recommendation number, CVE count, publicity, effort, or dependency.
This is neither a universal HDF5 rating nor a certification or conclusion about
an unassessed release or deployment.
\end{notebox}

The authoritative register contains fourteen formal findings: seven Rated and
seven Provisional. Their tiers comprise one Critical, four High, four
Moderate--High ranges, and five Moderate. Four additional observations remain
in the evidence-and-scoping backlog. These
counts describe the bounded scenarios in Appendix~\ref{app:findings-register};
they are not prevalence estimates and must not be aggregated into a score for
HDF5 as a technology.

\subsection{Reconciled Mitigation Priorities}

Table~\ref{tab:conclusion-priority-reconciliation} summarizes the only
report-wide ordering. The deployment census and assert-site source catalog are
due within 30 days of roadmap adoption; a finding-response clock starts when an
affected deployment or equivalent exposure is identified. Treatment removes or
controls the concrete scenario; \code{RM-1A} establishes scope and is not a
substitute for the operator-owned \code{RM-1F} containment. Prevention addresses
recurrence and may continue after the scoped finding is reassessed.

\begingroup
\small
\setlength{\tabcolsep}{3pt}
\renewcommand{\arraystretch}{1.15}

\setlength{\LTleft}{-0.9in}
\setlength{\LTright}{-0.9in}

\begin{longtable}{@{}>{\raggedright\arraybackslash}p{0.18\widetablewidth}>{\raggedright\arraybackslash}p{0.26\widetablewidth}>{\raggedright\arraybackslash}p{0.31\widetablewidth}>{\raggedright\arraybackslash}p{0.19\widetablewidth}@{}}
\caption{Conclusion-level reconciliation of mitigation priorities}
\label{tab:conclusion-priority-reconciliation}\\
\hline
\textbf{Decision lane} & \textbf{Records} & \textbf{Required outcome} &
\textbf{Boundary} \\
\hline
\endfirsthead

\caption[]{Conclusion-level reconciliation of mitigation priorities, continued}\\
\hline
\textbf{Decision lane} & \textbf{Records} & \textbf{Required outcome} &
\textbf{Boundary} \\
\hline
\endhead

\hline
\endfoot

\hline
\endlastfoot

Immediate---before the next untrusted job and within 24 hours of identification &
\code{APP-SEC-001}; any separately scoped service scenario that is rerated
Critical under the report-wide rubric &
Stop or sandbox attacker-selected external-resource resolution; remove
unnecessary worker authority, preserve evidence, and rotate exposed secrets
where indicated. &
Applies to an equivalent still-exposed deployment, not automatically to every
HDF5 service or to Hugging Face's present state. \\

Near-term---treatment within 30 days of affected-scope identification &
\code{CORE-SEC-001}, \code{CORE-TOOL-001}, \code{CORE-PRIV-001},
\code{CORE-SAFE-001}, \code{CORE-SAFE-002},
\code{APP-SEC-002}, \code{APP-SEC-003}, \code{APP-SEC-005} &
Use the program census to identify scope; patch, upgrade, isolate, or deploy a
tested compensating control for each affected deployment; minimize sensitive
plaintext metadata and verify the enforcing trust boundary. For the two
write-side safety findings the obligation is tested operational containment ---
capacity headroom and monitoring, flush and checkpoint discipline, a rehearsed
recovery procedure, and an independent copy of unregenerable data under
\code{RM-1F} --- not the library change or the \code{RM-1A} census. &
Census: 30 days after roadmap adoption. Treatment: 30 days after affected-scope
identification. Ranges route at their upper bound; dependencies do not extend
treatment. \\

Planned---within 90 days or the next supported release or migration window,
whichever is sooner &
\code{CORE-SEC-002}, \code{APP-SEC-004}, \code{IND-SAF-001},
\code{IND-SAF-002}, \code{LONG-SAF-001} &
Bound file-controlled work and allocation; verify exact reader/writer behavior;
inventory affected artifacts and filter dependencies; and provide tested
repair, migration, dependency-preservation, or controlled compatibility
handling. &
Expedite a scoped item when an affected production pipeline, archive, or
service is identified. \\

Evidence action---no assigned risk urgency &
\code{CORE-SEC-003}, \code{FMT-SAF-001}, \code{EXT-PRIV-001},
\code{SUP-SAFE-001} &
Perform bounded, reproducible studies sufficient to rate and route a supported
scenario or to document a negative result over the tested scope; catalog and
compare assertion-enabled and \code{-DNDEBUG} behavior for file-derived
assertions; adjudicate the identified normative dispute and measure any
dependent file or workflow consequence. &
NE and backlog status do not mean Low risk. The 30-day Phase~1 source-catalog
target and Phase~2 administrative targets are assurance schedules, not
risk-response categories. \\

Continuing prevention and sustainment &
Relevant formal findings and conditional evidence records &
Deliver strict processing and preflight controls, invariant-driven parsing,
always-active release-build checks, checked arithmetic, differential
conformance tests, governed extension loading, privacy-lifecycle controls,
provenance, release gates, and reassessment. &
Phases~2--4 may run in parallel. Structural work does not defer earlier
treatment or acquire an inherent-risk tier of its own. \\

\end{longtable}
\endgroup

This ordering also reconciles the chapters. The application chapter owns the
only Immediate route. The core and tool chapter combines Near-term treatment of
confirmed affected High and upper-bound-High scope with Planned resource controls and
longer-term class prevention; the assert-masking observation adds a 30-day
source-catalog action but no risk tier of its own; its development-tree
reproduction remains outside the rated population until a supported release
path is established. The core chapter additionally owns two Near-term write-side safety
containments and their longer structural treatments. The
independent-implementation chapter owns two Planned interoperability
treatments and one format-governance evidence action, and the long-term-access
chapter owns one Planned, horizon-qualified, filter-specific archival finding
that an earlier revision left unregistered. Extension privacy and supply-chain mismatch themes
remain evidence actions until a concrete affected population and harm are
established. The wrappers chapter creates no additional formal priority in this revision. Appendix~\ref{app:phased-roadmap} then
translates these lanes into accountable work packages without changing their
inherent risk.

\subsection{Significance of the 2026 Developments}

The Hugging Face incident and HDF5 issue \#6618 sharpen the application-boundary
conclusion. The incident demonstrates that a valid, documented external-storage
capability can become a repeatable disclosure mechanism when an automatic
service supplies filesystem authority on behalf of an untrusted file. That
evidence supports Immediate containment of equivalent deployments and the need
for unified, enforcement-coupled external-resource policy. It does not support
reclassifying the mechanism as a core parser vulnerability, treating a policy
proposal as an implemented control, or assigning the incident's Critical tier
to every HDF5 consumer
\cite{huggingface-agent-intrusion-2026,
hdfgroup-mechanism-not-vulnerability-2026,h5py-gh-issue-2891,
hdf5-gh-issue-6618}.

The PureHDF case sharpens the safety and interoperability conclusion in a
different way. Older-reader acceptance of contradictory chunk metadata did not
make the file conforming, while the separate sufficient-but-nonminimal case
shows that stricter validation can itself reject a permitted encoding. Together
the cases require a normative invariant decision, paired positive and negative
fixtures, exact supported-reader verification, an inventory of affected legacy
files, and deterministic repair or migration. The correct response is neither
permissiveness by default nor strictness without compatibility evidence
\cite{hdf5-gh-issue-6409,hdf5-gh-pr-6432,hdf5-gh-issue-6534,
purehdf-gh-issue-164}.

HDF5 Pickles issue~\#87 sharpens the core-parser conclusion. A validity
predicate derived from in-file bytes but enforced only inside
\code{assert(...)} is not a release control when \code{-DNDEBUG} removes it.
The issue's seven component families and cursory 50--100 estimate define an
urgent audit perimeter, not 50--100 proven vulnerabilities. This changes
sequence without manufacturing a rating: complete the catalog and paired
release/debug tests within 30 days of roadmap adoption, route and locally
repair a supported adverse path as soon as it is confirmed, and use the
results to order independently releasable Phase~3 migration tranches
~\cite{hdf5-pickles-issue-87}.

The perimeter is not entirely static. One source-reported case shows a
version-2 B-tree zero record size reaching an integer division whose only guard
is an assertion, producing \code{SIGFPE} through a public API path on an
\code{-DNDEBUG}, release-configured HDF5 2.3.0 development tree
~\cite{h5policy-case-v2btree-recordsize}. The same case record states that
32-bit behavior, affected supported releases, and the version range were not
measured. It therefore corroborates the mechanism and supplies a concrete
fixture, but it does not meet this report's supported-version condition for
promotion into \code{CORE-SEC-002}. That distinction demonstrates the rubric
working as designed: useful development-tree evidence can sharpen an audit
perimeter without acquiring a tier before the rated population is established.

The developments therefore reinforce the rubric rather than reorder it by
recency or publicity. The Hugging Face evidence supports a Critical
application-boundary scenario with an Immediate conditional response; the two
independent-implementation cases support Moderate safety scenarios with Planned
treatment; issue~\#6616 creates an administrative adjudication and evidence
action without a risk tier; and the assert-masking observation remains an
evidence action until a supported-release path and consequence are verified.
Together they show why the enforcing actor, affected scope, consequence, and
acceptance evidence must be explicit.

\subsection{Assurance, Closure, and Residual Risk}

This report does not establish that any proposed mitigation has been adopted or
that current residual risk has been reduced. Unless a record expressly states
otherwise, the assessment relied on public document and artifact examination
rather than independent reproduction. Current installed populations, exact fix
coverage across supported releases, and deployment-specific control
effectiveness remain incomplete. Every roadmap package is therefore
\emph{Proposed / implementation not verified} at publication, and no
residual-risk acceptance is recorded.

Priority becomes an auditable commitment only after the roadmap authority
adopts or modifies the plan, defines the census boundary, assigns named owners
and due dates, and preserves finding-by-scope response records. A concrete
finding may close after verified local treatment removes its scoped scenario
and the accountable authority records the residual-risk disposition; the
separate recurrence-prevention program remains open until its own acceptance
evidence passes. Every package requires independent verification and a
residual-risk disposition. For High and Critical scope, the verifier must also
be separate from the package owner, rollout approver, and risk acceptor.

The practical conclusion is consequently twofold. Known affected populations
should be treated on their response clocks without waiting for architectural
work, and HDF5's recurring classes should be addressed structurally without
pretending that local closure proves class elimination. Applied together, those
tracks convert the report from a catalog of concerns into a testable mitigation
program: constrain file-selected capabilities at authority boundaries, validate
file-derived state before trusted object construction, decide compatibility
against normative evidence, and govern privacy across the full data lifecycle.

\end{document}
